\section{Elementary estimates}\label{sec:elementary}
For a finite lattice domain $D$, write
\[
 G_D(x,y)=E_x\sum_{j<\tau_{D^c}}\mathbf1_{\{S_j=y\}},
 \qquad \tau_{D^c}=\inf\{j\ge0:S_j\notin D\}.
\]
The following estimates fix the normalization used throughout the paper.
\begin{lemma}[Elementary walk estimates]\label{lem:elementary-walk-estimates}
There are constants \(c,C\in(0,\infty)\) with the following properties.

\begin{enumerate}\item If \(\tau_R=\inf\{n:\lVert S_n\rVert_2\ge R\}\), then
   \[
       \sup_{\lVert x\rVert<R}E_x\tau_R\le C R^2,\qquad
       \sup_{\lVert x\rVert<R}P_x(\tau_R>mCR^2)\le 2^{-m}.
   \]
\item For all \(N,R\ge1\),
   \[
       P\left(\max_{k\le N}\lVert S_k\rVert_2\ge R\right)
       \le C\exp\{-cR^2/N\}.
   \]
\item Uniformly for \(1\ll r<\lVert x\rVert<R\),
   \[
   P_x(\tau_{\partial D(0,R)}<\tau_{\partial D(0,r)})
     =\frac{\log(\lVert x\rVert/r)+O(r^{-1})}{\log(R/r)}
   \]
   and the complementary inner-hitting probability is obtained by replacing
   the numerator by \(\log(R/\lVert x\rVert)+O(r^{-1})\).
\item For each fixed $\kappa\in(0,1)$, uniformly for $\|x\|\le(1-\kappa)R$ in
   \(D(0,R)\),
   \[
       G_{D(0,R)}(x,x)=g\log R+O_\kappa(1).
   \]
   Conditional on first hitting \(x\), the number of visits to \(x\) before
   leaving the disk is geometric with this mean.
\end{enumerate}
\end{lemma}

\begin{proof}
For (1), optional stopping of
\(\lVert S_{n\wedge\tau_R}\rVert_2^2-(n\wedge\tau_R)\), followed by the
one-step overshoot bound, gives the expectation estimate.  Markov's
inequality at time \(2CR^2\), iterated with the strong Markov property,
gives the geometric tail.  Part (2) follows by applying the usual
exponential martingale separately to both coordinates and to both signs.

For the annulus estimate, we use \cite[Lemma~2.1, (2.5)]{HLOZ2025}:
uniformly for $r<\|x\|<R$,
\[
 P_x(H_{\partial D(0,R)}<H_{\partial D(0,r)})
 =\frac{\log(\|x\|/r)+O(r^{-1})}{\log(R/r)}.
\]
Here \(H\) is the positive-time hitting time and the boundaries are the
digital outer vertex boundaries defined in that paper.  The later
directed-edge convention can differ from this vertex-boundary convention
only within a one-unit radial collar.  Applying the estimate with \(r,R\)
shifted by at most two units, equivalently repeating its potential-kernel
proof at the directed crossing time, changes the logarithmic terms by
\(O(r^{-1})\).  Hence the same estimate holds for the directed-edge rows
below.  The potential-kernel proof of that result also gives (4).  The
geometric assertion in (4) is the strong Markov
property at successive positive returns to \(x\).
\end{proof}
